% ============================================
% Глава 2. Проектирование системы sdCloud
% ============================================

\chapter{ПРОЕКТИРОВАНИЕ СИСТЕМЫ SDCLOUD}

\section{Общая архитектура платформы}

\subsection{Принципы проектирования}

Проектирование программной платформы \texttt{sdCloud} основано на совокупности архитектурных принципов, обеспечивающих масштабируемость, поддерживаемость и расширяемость системы:

\begin{itemize}
    \item \textbf{Разделение ответственности} (Separation of Concerns) --- чёткое разграничение функциональности между компонентами системы: frontend отвечает за представление и взаимодействие с пользователем, backend реализует бизнес-логику, а внешние сервисы предоставляют специализированные возможности (LLM, хранение данных).
    \item \textbf{Многоуровневая архитектура} (Layered Architecture) --- backend построен по принципу разделения на слои: слой API (контроллеры), слой сервисов (бизнес-логика), слой доступа к данным (репозитории) и слой моделей данных.
    \item \textbf{Асинхронная обработка} --- длительные операции (генерация моделей через LLM) выполняются асинхронно через очередь задач, что обеспечивает отзывчивость интерфейса и масштабируемость системы.
    \item \textbf{Модульность} --- система декомпозирована на функциональные модули с чётко определёнными интерфейсами, что упрощает тестирование и позволяет заменять реализации отдельных компонентов.
\end{itemize}

\subsection{Компонентная архитектура}

Платформа \texttt{sdCloud} состоит из следующих основных компонентов (рис.~\ref{fig:architecture}):

\begin{figure}[H]
    \centering
    \begin{tikzpicture}[
        node distance=0.8cm and 1.2cm,
        box/.style={draw, rounded corners=4pt, minimum width=3.2cm, minimum height=0.9cm, align=center, font=\small},
        layer/.style={draw, dashed, rounded corners=6pt, inner sep=8pt},
        arrow/.style={-{Stealth[scale=0.8]}, thick}
    ]
        % Client Layer
        \node[box, fill=blue!10] (client) {Веб-браузер\\(Next.js SPA)};

        % API Gateway Layer
        \node[box, fill=orange!10, below=of client] (api) {FastAPI Gateway\\REST API};

        % Service Layer
        \node[box, fill=green!10, below left=0.8cm and 0.3cm of api] (session) {Сервис\\сессий};
        \node[box, fill=green!10, below=of api] (gen) {Сервис\\генерации};
        \node[box, fill=green!10, below right=0.8cm and 0.3cm of api] (validate) {Сервис\\валидации};

        % Task Queue
        \node[box, fill=purple!10, below=1.5cm of gen] (taskiq) {Taskiq\\(очередь задач)};

        % Data Layer
        \node[box, fill=gray!15, below left=0.8cm and 0.2cm of taskiq] (postgres) {PostgreSQL\\(данные)};
        \node[box, fill=gray!15, below=of taskiq] (redis) {Redis\\(кэш, состояния)};

        % External
        \node[box, fill=red!10, right=2.5cm of api] (openrouter) {OpenRouter API\\(LLM)};
        \node[box, fill=red!10, below=of openrouter] (pysd) {PySD\\(симуляция)};

        % Arrows
        \draw[arrow] (client) -- (api);
        \draw[arrow] (api) -- (session);
        \draw[arrow] (api) -- (gen);
        \draw[arrow] (api) -- (validate);
        \draw[arrow] (gen) -- (taskiq);
        \draw[arrow] (taskiq) -- (postgres);
        \draw[arrow] (taskiq) -- (redis);
        \draw[arrow, dashed] (gen.east) -- ++(0.5,0) |- (openrouter.west);
        \draw[arrow, dashed] (taskiq.east) -- ++(0.5,0) |- (pysd.west);
        \draw[arrow] (session.south) -- ++(0,-0.4) -| (postgres.north);

    \end{tikzpicture}
    \caption{Архитектура платформы sdCloud}
    \label{fig:architecture}
\end{figure}

\begin{enumerate}
    \item \textbf{Веб-клиент} (Frontend) --- одностраничное приложение на базе Next.js, реализующее пользовательский интерфейс для создания сессий моделирования, ввода текстовых описаний, визуализации сгенерированных моделей (CLD, SFD, уравнения, графики симуляций) и управления результатами. Поддерживаются как авторизованный, так и гостевой режимы.

    \item \textbf{API Gateway} --- REST API на базе FastAPI, обеспечивающий точки входа для всех операций платформы, аутентификацию на основе JWT, валидацию входных данных и маршрутизацию запросов к внутренним сервисам.

    \item \textbf{Сервис сессий} --- управление сессиями моделирования, включающее создание, чтение, обновление и удаление сессий (CRUD-операции), а также управление сообщениями в рамках каждой сессии.

    \item \textbf{Сервис генерации} --- реализация конвейера генерации системно-динамических моделей, включающего взаимодействие с LLM через OpenRouter API и последовательное формирование переменных, CLD, уравнений, SFD, Python-кода и Vensim-экспорта.

    \item \textbf{Сервис валидации} --- проверка корректности графовых структур сгенерированных моделей с использованием алгоритмов на графах (библиотека NetworkX) и проверка согласованности между переменными, уравнениями и CLD.

    \item \textbf{Детерминированный генератор Python-кода} --- модуль, формирующий исполняемый Python-код модели по структурированным уравнениям без повторного обращения к LLM.

    \item \textbf{Симуляционный движок} --- PySD используется в качестве основного средства симуляции; при невозможности загрузки сгенерированного кода в PySD (например, из-за изменений в API версии 3.x) автоматически применяется fallback на SciPy \texttt{solve\_ivp}.

    \item \textbf{PostgreSQL} --- реляционная база данных для персистентного хранения пользователей, сессий, сообщений, моделей и результатов симуляций.

    \item \textbf{Redis} --- хранилище кэша и состояний, используемое для кэширования часто запрашиваемых данных.

    \item \textbf{OpenRouter API} --- внешний сервис-агрегатор, предоставляющий унифицированный доступ к различным LLM через единый REST API; взаимодействие реализовано через адаптер на базе \texttt{litellm} и OpenAI SDK.

    \item \textbf{Гостевой режим} --- набор stateless API-эндпоинтов, позволяющих анонимным пользователям синхронно генерировать модели и запускать симуляции без записи данных в PostgreSQL.
\end{enumerate}

\subsection{Взаимодействие компонентов}

Основной сценарий взаимодействия компонентов при генерации модели реализуется следующим образом:

\begin{enumerate}
    \item Пользователь вводит текстовое описание системы через веб-интерфейс.
    \item Frontend отправляет HTTP-запрос к API Gateway (авторизованный или гостевой).
    \item Для авторизованных пользователей API создаёт фоновую \texttt{asyncio}-задачу генерации; для гостей генерация выполняется синхронно в рамках одного запроса.
    \item Воркер генерации последовательно выполняет этапы конвейера, обращаясь к OpenRouter API для извлечения переменных, генерации CLD и уравнений, и строит SFD на основе уравнений.
    \item После генерации CLD и SFD выполняется валидация графовых структур и согласованности модели.
    \item Детерминированный генератор формирует Python-код модели; отдельный модуль экспортирует модель в формат Vensim .mdl.
    \item Сгенерированный Python-код передаётся в PySD; при ошибках загрузки автоматически используется SciPy \texttt{solve\_ivp} fallback.
    \item Результаты сохраняются в PostgreSQL (для авторизованных пользователей) или возвращаются напрямую в ответе (гостевой режим).
    \item Frontend отображает CLD, SFD, уравнения, Python/Vensim код и результаты симуляции.
\end{enumerate}


\section{Backend: реализация серверной части}

\subsection{Фреймворк FastAPI}

Серверная часть платформы реализована на базе фреймворка \textbf{FastAPI} \cite{fastapi2024} --- современного веб-фреймворка для Python, основанного на типизации (type hints) и стандарте ASGI. FastAPI обеспечивает следующие преимущества для проекта:

\begin{itemize}
    \item \textbf{Автоматическая валидация данных} на основе аннотаций типов Python и библиотеки Pydantic, что сокращает объём шаблонного кода и повышает надёжность.
    \item \textbf{Автоматическая генерация документации} API в форматах OpenAPI и Swagger UI, упрощающая разработку и тестирование.
    \item \textbf{Нативная поддержка асинхронного программирования} (\texttt{async}/\texttt{await}), что критично для эффективной работы с внешними API и базами данных.
    \item \textbf{Высокая производительность} --- FastAPI демонстрирует одни из лучших показателей пропускной способности среди Python-фреймворков.
\end{itemize}

\subsection{ORM SQLAlchemy и миграции Alembic}

Для взаимодействия с реляционной базой данных используется \textbf{SQLAlchemy} 2.0 \cite{sqlalchemy2024} --- наиболее зрелый и функциональный ORM-фреймворк для Python. SQLAlchemy 2.0 предоставляет:

\begin{itemize}
    \item Декларативное описание моделей данных с поддержкой типизации.
    \item Мощный язык запросов, генерирующий оптимальный SQL-код.
    \item Поддержку асинхронных сессий (\texttt{AsyncSession}) для неблокирующего доступа к базе данных.
    \item Встроенную систему миграций через интеграцию с Alembic.
\end{itemize}

\textbf{Alembic} используется для управления схемой базы данных. Система миграций обеспечивает версионирование структуры БД, возможность наката и отката изменений, а также согласованность схемы между различными окружениями (разработка, тестирование, продакшн).

\subsection{База данных PostgreSQL}

\textbf{PostgreSQL} \cite{postgresql2024} выбран в качестве основной системы управления базами данных благодаря следующим характеристикам:

\begin{itemize}
    \item Полная поддержка стандарта SQL и расширенные возможности (JSON/JSONB, массивы, пользовательские типы).
    \item Высокая надёжность и производительность.
    \item Поддержка асинхронных подключений через \texttt{asyncpg}.
    \item Расширяемость и богатая экосистема.
\end{itemize}

В контексте проекта \texttt{sdCloud} PostgreSQL хранит информацию о пользователях, сессиях моделирования, сообщениях, сгенерированных моделях (CLD, SFD, уравнения), результатах симуляций и метаданных.

\subsection{Redis}

\textbf{Redis} \cite{redis2024} используется в двух качествах:

\begin{enumerate}
    \item \textbf{Хранилище состояний задач} --- отслеживание статуса асинхронных задач генерации (в очереди, в процессе, завершено, ошибка).
    \item \textbf{Кэш} --- кэширование часто запрашиваемых данных (например, списка сессий пользователя) для снижения нагрузки на PostgreSQL.
\end{enumerate}

Асинхронный клиент \texttt{redis-py} обеспечивает неблокирующее взаимодействие с хранилищем в контексте async-приложения FastAPI.

\subsection{Асинхронная очередь задач Taskiq}

\textbf{Taskiq} \cite{taskiq2024} --- современная асинхронная очередь задач для Python, выбранная в качестве альтернативы Celery. Ключевые преимущества Taskiq:

\begin{itemize}
    \item Нативная поддержка \texttt{async}/\texttt{await} без необходимости в синхронных обёртках.
    \item Интеграция с FastAPI через middleware.
    \item Гибкая система брокеров (Redis, RabbitMQ и др.).
    \item Поддержка периодических задач (cron).
    \item Типизированные задачи с автоматической сериализацией.
\end{itemize}

В проекте \texttt{sdCloud} Taskiq используется для асинхронного выполнения конвейера генерации моделей. Каждый запрос на генерацию создаёт задачу, которая выполняется воркером независимо от HTTP-запроса пользователя.


\section{Frontend: клиентская часть}

\subsection{Фреймворк Next.js}

Клиентская часть реализована на базе \textbf{Next.js} \cite{nextjs2024} --- React-фреймворка, предоставляющего серверный рендеринг (SSR), статическую генерацию (SSG) и клиентскую навигацию. В контексте проекта \texttt{sdCloud} Next.js используется в режиме одностраничного приложения (SPA) с клиентской маршрутизацией и динамической загрузкой данных.

Выбор Next.js обусловлен следующими факторами:

\begin{itemize}
    \item Развитая экосистема инструментов и библиотек для React.
    \item Оптимизации производительности (код-сплиттинг, предзагрузка страниц).
    \item Гибкость в выборе стратегии рендеринга.
    \item Поддержка API Routes для серверных функций.
\end{itemize}

\subsection{Методология Feature-Sliced Design}

Архитектура frontend-приложения построена в соответствии с методологией \textbf{Feature-Sliced Design} (FSD) \cite{fsd2024} --- современным подходом к организации кода frontend-приложений, ориентированным на разделение по функциональным возможностям (фичам).

FSD определяет иерархическую структуру слоёв приложения:

\begin{table}[H]
    \centering
    \caption{Слои архитектуры FSD в проекте sdCloud}
    \label{tab:fsd-layers}
    \begin{tabular}{|p{3cm}|p{5cm}|p{5.5cm}|}
        \hline
        \textbf{Слой} & \textbf{Назначение} & \textbf{Примеры компонентов} \\
        \hline
        app & Инициализация приложения, провайдеры, глобальные стили & \texttt{providers/, router/, store/} \\
        \hline
        pages & Страницы приложения, композиция виджетов & \texttt{HomePage/, SessionPage/, ModelsPage/} \\
        \hline
        widgets & Самостоятельные блоки интерфейса & \texttt{ChatWidget/, ModelViewer/, Sidebar/} \\
        \hline
        features & Функциональные модули, реализующие бизнес-логику & \texttt{auth/, modelGenerate/, sessionCreate/} \\
        \hline
        entities & Бизнес-сущности предметной области & \texttt{session/, model/, user/} \\
        \hline
        shared & Переиспользуемые компоненты, утилиты, конфигурация & \texttt{api/, ui/, lib/, config/} \\
        \hline
    \end{tabular}
\end{table}

Ключевым правилом FSD является направленность зависимостей: компоненты нижнего слоя не могут зависеть от компонентов верхнего слоя. Это обеспечивает слабую связанность и высокую когезию модулей.

\subsection{Компонентная база shadcn/ui}

Для ускорения разработки интерфейса используется библиотека \textbf{shadcn/ui} \cite{shadcn2024} --- коллекция переиспользуемых React-компонентов, основанная на Radix UI primitives и стилизованная с помощью Tailwind CSS. Компоненты shadcn/ui копируются непосредственно в кодовую базу проекта, что обеспечивает полный контроль над их внешним видом и поведением.

\subsection{Управление состоянием и запросами}

Для управления серверным состоянием используется \textbf{TanStack Query} (React Query) \cite{tanstack2024} --- библиотека, предоставляющая декларативный подход к загрузке, кэшированию и обновлению данных с сервера. TanStack Query автоматизирует обработку состояний загрузки, ошибок и инвалидацию кэша.

Для глобального клиентского состояния используется библиотека \textbf{Zustand} --- лёгкое решение для управления состоянием с минимальным шаблонным кодом.


\section{Модуль генерации моделей}

\subsection{Конвейер генерации}

Модуль генерации реализует многоэтапный конвейер преобразования текстового описания системы в исполняемую системно-динамическую модель. Конвейер состоит из следующих последовательных этапов (рис.~\ref{fig:pipeline}):

\begin{figure}[H]
    \centering
    \begin{tikzpicture}[
        node distance=0.4cm,
        stage/.style={draw, rounded corners=4pt, minimum width=2.4cm, minimum height=1cm, align=center, font=\small},
        arrow/.style={-{Stealth[scale=0.8]}, thick}
    ]
        \node[stage, fill=blue!10] (text) {Текстовое\\описание};
        \node[stage, fill=green!10, right=0.6cm of text] (vars) {Переменные\\(LLM)};
        \node[stage, fill=green!10, right=0.6cm of vars] (cld) {CLD\\(LLM)};
        \node[stage, fill=green!10, right=0.6cm of cld] (eq) {Уравнения\\(LLM)};
        \node[stage, fill=green!10, right=0.6cm of eq] (sfd) {SFD\\(LLM/правила)};
        \node[stage, fill=yellow!15, right=0.6cm of sfd] (py) {Python\\(генератор)};
        \node[stage, fill=orange!10, right=0.6cm of py] (sim) {Симуляция\\(PySD/SciPy)};

        \draw[arrow] (text) -- (vars);
        \draw[arrow] (vars) -- (cld);
        \draw[arrow] (cld) -- (eq);
        \draw[arrow] (eq) -- (sfd);
        \draw[arrow] (sfd) -- (py);
        \draw[arrow] (py) -- (sim);
    \end{tikzpicture}
    \caption{Конвейер генерации системно-динамической модели (актуальная версия)}
    \label{fig:pipeline}
\end{figure}

\begin{enumerate}
    \item \textbf{Этап 1: Извлечение переменных.} На вход подаётся текстовое описание системы. LLM извлекает ключевые переменные с указанием типа (stock, flow, auxiliary, constant), единиц измерения, начальных значений и описания.

    \item \textbf{Этап 2: Генерация CLD.} На основе списка переменных LLM строит ориентированный граф причинно-следственных связей $G = (V, E)$, где $V$ --- множество переменных, $E \subseteq V \times V$ --- множество ориентированных рёбер с пометками полярности $(+, -)$. Граф валидируется на связность, наличие обратных связей и отсутствие изолированных вершин.

    \item \textbf{Этап 3: Генерация уравнений.} На основе переменных и CLD формируется система уравнений в структурированном JSON-формате. Каждое уравнение содержит левую часть (\texttt{lhs}), правую часть (\texttt{rhs}) и тип переменной. Уровни описываются через вызов \texttt{integ(чистый поток, начальное\_значение)}.

    \item \textbf{Этап 3b: Постпроцессинг уравнений.} Автоматически исправляются типичные ошибки LLM: константы и уровни с нулевыми или пустыми значениями заменяются на эвристически подобранные ненулевые defaults (например, 0.05 для коэффициентов, 100.0 для начальных запасов).

    \item \textbf{Этап 4: Проверка согласованности.} Проверяется соответствие между переменными, уравнениями и узлами CLD: каждая переменная имеет уравнение, каждое уравнение относится к известной переменной, все узлы CLD присутствуют в списке переменных.

    \item \textbf{Этап 5: Генерация SFD.} На основе переменных и уравнений строится диаграмма потоков и уровней. LLM формирует позиции узлов и связи, после чего выполняется валидация структуры SFD (наличие потоков для уровней, отсутствие изолированных элементов).

    \item \textbf{Этап 6: Детерминированная генерация Python-кода.} По структурированным уравнениям формируется исполняемый Python-код, совместимый с PySD. Все bare-ссылки на переменные преобразуются в вызовы функций от \texttt{time}, а функции \texttt{MAX}/\texttt{MIN}/\texttt{ABS} нормализуются к Python-встроенным \texttt{max}/\texttt{min}/\texttt{abs}. Данный этап не требует обращения к LLM, что делает генерацию воспроизводимой.

    \item \textbf{Этап 7: Генерация Vensim-кода.} Уравнения экспортируются в текстовый формат Vensim .mdl с преобразованием Python-функций (\texttt{np.exp}, \texttt{max} и др.) в Vensim-синтаксис.

    \item \textbf{Этап 8: Симуляция.} Python-код модели загружается в PySD; при ошибках, связанных с импортом \texttt{integ} или несовместимостью API, автоматически используется SciPy \texttt{solve\_ivp} fallback. Результаты симуляции (временные ряды) возвращаются пользователю.
\end{enumerate}

\subsection{Промпт-инжиниринг}

Качество генерации LLM критически зависит от структуры промптов. В проекте \texttt{sdCloud} используются \textbf{few-shot промпты}, включающие примеры входных данных и ожидаемых ответов для каждого этапа конвейера.

Структура промпта для генерации CLD:

\begin{lstlisting}[caption={Структура промпта для генерации CLD}]
SYSTEM: Ты --- эксперт по системной динамике. Твоя задача ---
построить диаграмму причинно-следственных связей (CLD) на основе
текстового описания системы.

Формат ответа: JSON со следующей структурой:
{
  "variables": ["var1", "var2", ...],
  "edges": [
    {"from": "var1", "to": "var2", "polarity": "+"},
    ...
  ],
  "loops": [
    {"type": "R", "variables": ["var1", "var2", "var3"]},
    ...
  ]
}

Пример:
Описание: "Популяция кроликов растет при наличии пищи.
Чем больше популяция, тем больше потребление пищи."

Ответ:
{
  "variables": ["Популяция кроликов", "Пищевые ресурсы",
               "Потребление пищи"],
  "edges": [
    {"from": "Пищевые ресурсы", "to": "Популяция кроликов",
     "polarity": "+"},
    {"from": "Популяция кроликов", "to": "Потребление пищи",
     "polarity": "+"},
    {"from": "Потребление пищи", "to": "Пищевые ресурсы",
     "polarity": "-"}
  ],
  "loops": [
    {"type": "R", "variables": ["Популяция кроликов",
     "Потребление пищи", "Пищевые ресурсы"]},
    {"type": "B", "variables": ["Популяция кроликов",
     "Потребление пищи", "Пищевые ресурсы"]}
  ]
}

Теперь построй CLD для следующего описания:
{user_description}
\end{lstlisting}

\subsection{Детерминированный генератор Python-кода}

Важным архитектурным решением, отличающим текущую версию платформы от ранних прототипов, является отказ от генерации исполняемого Python-кода непосредственно LLM. Вместо этого используется детерминированный генератор \texttt{python\_generator.py}, который строит код по структурированным уравнениям:

\begin{itemize}
    \item Каждая переменная превращается в функцию \texttt{def имя(time): return выражение}.
    \item Все bare-ссылки на переменные внутри правых частей автоматически преобразуются в вызовы функций, что устраняет проблему пропущенных скобок, характерную для LLM-генерации.
    \item Поддерживается нормализация Vensim-функций \texttt{MAX}, \texttt{MIN}, \texttt{ABS} в Python-встроенные аналоги.
    \item Результат воспроизводим: один и тот же набор уравнений всегда порождает идентичный код.
\end{itemize}

Этот подход существенно повышает надёжность симуляции, поскольку исключает целый класс ошибок, связанных с галлюцинациями LLM при генерации кода (неправильные импорты, неопределённые переменные, синтаксические ошибки).


\subsection{Fallback-симуляция через SciPy}

В процессе эксплуатации платформы выявлена несовместимость сгенерированного Python-кода (использующего \texttt{from pysd.py\_backend.functions import integ}) с актуальной версией PySD 3.14.x, в которой символ \texttt{integ} недоступен для импорта. Для сохранения работоспособности системы разработан fallback-модуль \texttt{scipy\_runner.py}:

\begin{itemize}
    \item Парсит сгенерированный Python-код через модуль \texttt{ast}.
    \item Выделяет уровни (stock) по наличию вызова \texttt{integ(поток, начальное\_значение)}.
    \item Формирует функцию \texttt{derivatives(t, state)}, вычисляющую чистые потоки для каждого уровня.
    \item Интегрирует систему ОДУ методами SciPy (RK45 по умолчанию, доступны LSODA, BDF, DOP853 и др.).
    \item Поддерживает переопределение параметров и начальных значений уровней.
\end{itemize}

Таким образом, симуляция остаётся работоспособной даже при отказе PySD, что повышает отказоустойчивость платформы.


\subsection{Форматы промежуточных представлений}

Промежуточные результаты каждого этапа конвейера сериализуются в формате JSON, обеспечивая прозрачность процесса и возможность отладки:

\textbf{Формат CLD} (этап 2):
\begin{lstlisting}[caption={Формат представления CLD}]
{
  "variables": ["Population", "Births", "Deaths",
                "Food", "Consumption"],
  "edges": [
    {"from": "Population", "to": "Births", "polarity": "+"},
    {"from": "Population", "to": "Deaths", "polarity": "+"},
    {"from": "Food", "to": "Births", "polarity": "+"},
    {"from": "Food", "to": "Consumption", "polarity": "+"},
    {"from": "Consumption", "to": "Food", "polarity": "-"}
  ],
  "loops": [
    {"type": "R", "variables": ["Population", "Births"]},
    {"type": "B", "variables": ["Population", "Deaths"]}
  ]
}
\end{lstlisting}

\textbf{Формат SFD} (этап 5):
\begin{lstlisting}[caption={Формат представления SFD}]
{
  "stocks": [
    {"name": "Population", "initial": 100,
     "inflows": ["Births"], "outflows": ["Deaths"]}
  ],
  "flows": [
    {"name": "Births", "formula": "Population * birth_rate"},
    {"name": "Deaths", "formula": "Population * death_rate"}
  ],
  "auxiliaries": [
    {"name": "birth_rate", "value": 0.05},
    {"name": "death_rate", "value": 0.02}
  ]
}
\end{lstlisting}

\textbf{Формат Vensim MDL} (этап 7):
\begin{lstlisting}[caption={Пример сгенерированного кода MDL}]
Population = INTEG(Births - Deaths, 100)
~ people
~

Births = Population * birth_rate
~ people/year
~

Deaths = Population * death_rate
~ people/year
~

birth_rate = 0.05
~ 1/year
~

death_rate = 0.02
~ 1/year
~
\end{lstlisting}


\section{Модуль валидации графов}

\subsection{Задачи валидации}

Модуль валидации графов выполняет проверку корректности топологии сгенерированных CLD и SFD перед переходом к генерации уравнений и симуляции. Валидация направлена на выявление типичных ошибок генерации LLM:

\begin{itemize}
    \item Переменные без входящих или исходящих связей (изолированные вершины).
    \item Некорректная классификация переменных (отсутствие уровней в SFD).
    \item Нарушение баланса потоков для уровней.
    \item Отсутствие замкнутых петель обратной связи.
    \item Ссылки на несуществующие переменные в уравнениях.
\end{itemize}

\subsection{Алгоритмы валидации}

Валидация реализована с использованием библиотеки \textbf{NetworkX} \cite{networkx2008} --- зрелой Python-библиотеки для создания, манипулирования и анализа графов.

\textbf{Проверка связности графа:}

Ориентированный граф CLD должен быть \emph{слабо связным} (weakly connected), то есть его неориентированная версия должна быть связным графом. Это условие гарантирует отсутствие изолированных подсистем:

\begin{algorithm}[H]
    \caption{Проверка слабой связности CLD}
    \begin{algorithmic}[1]
        \Require Ориентированный граф $G = (V, E)$
        \Ensure \text{True}, если граф слабо связен; \text{False} иначе
        \State $G_{undirected} \gets$ \Call{ToUndirected}{$G$}
        \State $components \gets$ \Call{ConnectedComponents}{$G_{undirected}$}
        \If{$|components| = 1$}
            \State \Return \text{True}
        \Else
            \State \Return \text{False}
        \EndIf
    \end{algorithmic}
\end{algorithm}

\textbf{Обнаружение петель обратной связи:}

Системно-динамическая модель должна содержать хотя бы одну петлю обратной связи. Для обнаружения петель используется алгоритм поиска элементарных циклов в ориентированном графе (алгоритм Джонсона \cite{johnson1975finding}):

\begin{algorithm}[H]
    \caption{Обнаружение петель обратной связи}
    \begin{algorithmic}[1]
        \Require Ориентированный граф $G = (V, E)$
        \Ensure Список циклов $C$
        \State $C \gets$ \Call{SimpleCycles}{$G$} \Comment{алгоритм Джонсона}
        \State $feedback\_loops \gets \emptyset$
        \For{$cycle$ \textbf{in} $C$}
            \State $sign \gets$ \Call{ComputeLoopSign}{$cycle, E$}
            \State \Call{Append}{$feedback\_loops, (cycle, sign)$}
        \EndFor
        \State \Return $feedback\_loops$
    \end{algorithmic}
\end{algorithm}

Знак петли определяется произведением полярностей её рёбер: чётное число отрицательных связей соответствует усиливающей петле ($R$), нечётное --- балансирующей ($B$).

\textbf{Проверка структуры SFD:}

Для SFD выполняется проверка корректности структурных элементов:

\begin{itemize}
    \item Каждый уровень имеет хотя бы один входной или выходной поток.
    \item Каждый поток связан ровно с одним уровнем (источником или стоком).
    \item Все переменные, упомянутые в формулах потоков, определены в модели.
\end{itemize}

\subsection{Результаты валидации}

Результаты проверки графа возвращаются в структурированном формате:

\begin{lstlisting}[caption={Формат результата валидации}]
{
  "valid": false,
  "errors": [
    {"type": "isolated_variable", "variable": "ExternalFactor"},
    {"type": "unbalanced_stock",
     "stock": "Inventory", "missing": "outflow"}
  ],
  "warnings": [
    {"type": "no_reinforcing_loop",
     "message": "Model has no reinforcing loops"}
  ],
  "stats": {
    "variables": 8,
    "edges": 12,
    "feedback_loops": 3,
    "reinforcing": 1,
    "balancing": 2
  }
}
\end{lstlisting}

При обнаружении ошибок модуль валидации формирует отчёт, который может быть использован для корректирующей генерации (iterative refinement) --- повторного запроса к LLM с указанием выявленных проблем.


\section{Интеграция с OpenRouter}

\subsection{Архитектура интеграции}

\textbf{OpenRouter} --- сервис-агрегатор, предоставляющий унифицированный REST API для доступа к различным LLM (GPT-4, Claude, LLaMA, Mistral, Gemma и др.). Взаимодействие с OpenRouter в платформе \texttt{sdCloud} реализовано через двухуровневый адаптер:

\begin{itemize}
    \item \textbf{UniversalAI} --- синхронный клиент, использующий библиотеку \texttt{litellm} и официальный OpenAI SDK для вызова чат-комплишенов, генерации изображений, аудио, видео и транскрипции.
    \item \textbf{OpenRouterLLM} --- асинхронная обёртка над \texttt{UniversalAI}, выполняющая блокирующие вызовы в отдельном потоке (\texttt{asyncio.to\_thread}) с заданным таймаутом, что позволяет использовать клиент в FastAPI без блокировки event loop.
\end{itemize}

Использование OpenRouter в качестве прослойки между платформой \texttt{sdCloud} и провайдерами LLM обеспечивает следующие преимущества:

\begin{itemize}
    \item \textbf{Унификация интерфейса} --- единый формат запросов и ответов для всех моделей через совместимый с OpenAI API.
    \item \textbf{Гибкость выбора модели} --- возможность переключения между различными LLM изменением только переменной окружения \texttt{OPENROUTER\_MODEL}.
    \item \textbf{Доступ к бесплатным моделям} --- возможность использования free-tier моделей (например, \texttt{openrouter/google/gemma-4-31b-it:free}) для прототипирования и тестирования.
    \item \textbf{Расширенные модальности} --- потенциал для генерации сопроводительных изображений, аудио- и видеоматериалов к моделям.
\end{itemize}

\subsection{Формат запросов}

Взаимодействие с OpenRouter API реализовано через HTTP-запросы к эндпоинту \texttt{/chat/completions}, совместимому с форматом OpenAI API. Пример запроса, формируемого через \texttt{litellm}:

\begin{lstlisting}[caption={Пример запроса к OpenRouter API}]
POST https://openrouter.ai/api/v1/chat/completions
Content-Type: application/json
Authorization: Bearer {OPENROUTER_API_KEY}

{
  "model": "openrouter/google/gemma-4-31b-it:free",
  "messages": [
    {"role": "system", "content": "system_prompt"},
    {"role": "user", "content": "user_description"}
  ],
  "temperature": 0.3,
  "max_tokens": 4096
}
\end{lstlisting}

Параметр \texttt{temperature} установлен на уровне 0.3--0.4 для снижения стохастичности генерации и повышения воспроизводимости результатов. Для извлечения переменных, CLD, уравнений и SFD ответ модели извлекается из markdown-блока \texttt{```json} и парсится как JSON; при неудаче выполняется fallback-поиск массива или объекта в тексте.

\subsection{Обработка ответов и ошибок}

Ответ от OpenRouter API содержит сгенерированный текст в поле \texttt{choices[0].message.content}. Модуль интеграции выполняет следующие операции:

\begin{enumerate}
    \item Извлечение JSON из markdown-блоков или прямой поиск JSON-объекта/массива в тексте.
    \item Валидация структуры ответа (наличие обязательных полей, корректность типов).
    \item Обработка ошибок API (таймауты, лимиты запросов, недоступность модели, ошибки аутентификации).
    \item Логирование запросов и ответов для отладки.
\end{enumerate}

При возникновении ошибок генерация конкретного артефакта не прерывает весь конвейер: пустые списки переменных, уравнений или рёбер заменяются на безопасные defaults, а информация об ошибке фиксируется в логе генерации.


\section{Гостевой режим работы}

В текущей версии платформы реализован гостевой режим, позволяющий анонимным пользователям знакомиться с возможностями \texttt{sdCloud} без создания учётной записи.

\subsection{Принципы гостевого режима}

\begin{itemize}
    \item \textbf{Statelessness} --- сервер не сохраняет данные гостевых пользователей в PostgreSQL.
    \item \textbf{История в браузере} --- полный контекст диалога и сгенерированные модели хранятся в памяти браузера (Zustand store).
    \item \textbf{Синхронная генерация} --- в отличие от авторизованного режима, гостевая генерация выполняется синхронно в рамках одного HTTP-запроса, что упрощает клиентскую логику.
    \item \textbf{Полный набор артефактов} --- гость получает переменные, CLD, уравнения, SFD, Python-код, Vensim-код и может запустить симуляцию.
\end{itemize}

\subsection{API гостевого режима}

Гостевые эндпоинты расположены по префиксу \texttt{/api/v1/guest}:

\begin{itemize}
    \item \texttt{POST /api/v1/guest/chat} --- отправка сообщения с синхронной генерацией модели и ответом ассистента.
    \item \texttt{POST /api/v1/guest/simulate} --- запуск симуляции по предоставленному Python- или Vensim-коду.
\end{itemize}

На frontend гостевой режим поддерживается отдельным слоем сущностей \texttt{entities/guest} и виджетом \texttt{chat-layout}, который корректно обрабатывает отсутствие авторизации и сохраняет историю в Zustand.
